% Generated draft: requires mathematical and source review.
\documentclass[11pt]{article}
\usepackage[margin=25mm]{geometry}
\usepackage{fontspec}
\setmainfont{Latin Modern Roman}
\usepackage{amsmath,amssymb,mathtools,unicode-math}
\setmathfont{Latin Modern Math}
\usepackage{enumitem,needspace,xurl,hyperref}
\setlength{\parindent}{0pt}
\setlength{\parskip}{5pt}
\setlength{\emergencystretch}{3em}
\newcommand{\N}{\mathbb{N}}
\newcommand{\Z}{\mathbb{Z}}
\newcommand{\Q}{\mathbb{Q}}
\newcommand{\R}{\mathbb{R}}
\newcommand{\C}{\mathbb{C}}
\newcommand{\F}{\mathbb{F}}
\newcommand{\A}{\mathbb{A}}
\newcommand{\PP}{\mathbb{P}}
\newcommand{\E}{\mathbb{E}}
\newcommand{\eps}{\varepsilon}
\newcommand{\dd}{\,\mathrm{d}}
\begin{document}
\section*{UnsolvedMath ID: 20002727}
\textbf{Source problem code:} catalogue entry
\subsection{Definitions and mathematical statement}
D. Bessis and C. Kriloff ask what matrix mutation does to noncrystallographic finite-type Cartan data and whether there are recurrences. The source's statement that no cluster-algebra construction exists is historical context, not a current assessment. The Reading--Speyer remark reports approximate returns close to, but bounded away from, the initial data without specifying a metric, mutation word, or whether matrices or variables are meant.

A Coxeter Cartan matrix, with diagonal $2$, is not itself a zero-diagonal skew-symmetrizable exchange matrix. The aid first chooses an orientation and real exchange-matrix normalization. For a real exchange matrix $B$, write $[a]_+=\max(a,0)$ and use
\[
 (\mu_kB)_{ij}=\begin{cases}
 -b_{ij},&i=k\text{ or }j=k,\\
 b_{ij}+[b_{ik}]_+[b_{kj}]_+-[-b_{ik}]_+[-b_{kj}]_+,&i,j\ne k.
 \end{cases}
\]
These real matrix mutations are well-defined independently of whether ordinary rational cluster-variable mutations are available.

For every $p,q>0$, let $B_{p,q}=\left(\begin{smallmatrix}0&p\\-q&0\end{smallmatrix}\right)$. The reported exact rank-two statement is
\[
 \mu_1(B_{p,q})=\mu_2(B_{p,q})=-B_{p,q},\qquad
 \operatorname{Mut}(B_{p,q})=\{B_{p,q},-B_{p,q}\}.
\]
Thus the matrix returns exactly after any two mutations. If $pq=p'q'$, positive diagonal conjugacy relates the two normalizations, and
\[
 \mu_k(DBD^{-1})=D\mu_k(B)D^{-1}\quad(D\text{ positive diagonal}).
\]
For dihedral Coxeter data $I_2(m)$ the conventional invariant is $pq=4\cos^2(\pi/m)$, with symmetric choice $p=q=2\cos(\pi/m)$. The rank-two matrix recurrence therefore does not depend on crystallographic integrality.

Variable dynamics is different. For instance the associated real-power relations on positive variables include $x_1'=(1+x_2^q)/x_1$ and $x_2'=(1+x_1^p)/x_2$. If a relevant exponent is nonintegral, that mutation is not a rational self-map of the ordinary field $\R(x_1,x_2)$. Real-power dynamics can still be defined on a positive domain but need not have the matrix's two-step period. Positive diagonal matrix conjugacy does not automatically identify these variable dynamics. Any proposed approximate rank-two recurrence must therefore specify the extra dynamical object, not be attributed to this exact two-state matrix orbit.

The aid also gives a higher-rank symbolic illustration. With $\varphi=(1+\sqrt5)/2$, the oriented $H_3$ example
\[
 B=\begin{pmatrix}0&\varphi&0\\-\varphi&0&1\\0&-1&0\end{pmatrix}
 \quad\text{has}\quad
 \mu_2B=\begin{pmatrix}0&-\varphi&\varphi\\\varphi&0&-1\\-\varphi&1&0\end{pmatrix}.
\]
A new edge appears, so higher-rank mutation cannot be reduced to the rank-two sign-flip observation. This single calculation neither classifies the full higher-rank mutation classes nor supplies a general recurrence theorem.
\subsection{Short English statement}
The source must distinguish Cartan data, real exchange matrices, and cluster variables. Rank-two real matrix mutation has an exact two-state orbit, while nonintegral real-power variable updates are not ordinary rational cluster mutations. Higher-rank recurrence questions remain beyond that observation.
\subsection{Sources}
\begin{itemize}
\item Local UnsolvedMath registry, permanent ID 20002727, source code catalogue entry; created and updated 14 August 2026. Canonical \texttt{aim-probability-notes.json}, zero-based \texttt{notes[168]}, tag \texttt{problem}.
\item AIM ``Braid groups, clusters and free probability,'' extracted heading ``Problem 6.'' followed by ``4,'' attributed to D. Bessis and C. Kriloff: \url{https://aimath.org/WWN/braidgroups/braidgroups.pdf}. The terminal 13 is pagination. Statement status is reconstructed\_unverified. The original remark names N. Reading and D. Speyer for early approximate-recurrence calculations; literature field is empty.
\item Research attempt 1: \texttt{partial\_result}, problem status at run \texttt{partially\_solved}, rigor label \texttt{proved}, novelty confidence low. No complete authoritative bibliographic references, higher-rank classification, or quantitative approximate-recurrence statement are supplied.
\item Ulam AI / UnsolvedMath Contributors, CC BY 4.0. Machine-generated research aid, not peer reviewed; independently verify before citation.
\end{itemize}
\end{document}
